\section{The project}
\label{sec:intro}

\subsection{The problem}

People who learn technical subjects on their own accumulate notes faster than they accumulate
understanding. The notes end up scattered across a chat with themselves, a folder of screenshots, twelve
open browser tabs and three note-taking apps, and the connective tissue --- \emph{this depends on that},
\emph{I already know this part}, \emph{that term means something specific here} --- lives only in memory
and decays. Six months later the notes are still there and the thread through them is gone.

Two things are missing from the usual note app. The first is \textbf{structure that survives}: a way to
say that \emph{Raft} belongs under \emph{Consensus} under \emph{Distributed Systems}, that it relates to
\emph{Leader Election}, and that it is currently half-read rather than finished. The second is
\textbf{help with the part that is actually hard}: not writing the notes down, but working out what to
learn next, in what order, given what is already understood.

\subsection{Learning Mocha}

Learning Mocha is a native Android app that treats a person's own learning as a small private
encyclopaedia they own outright:

\begin{itemize}
  \item a \textbf{personal Wikipedia} --- Markdown articles organised into branches, folders and
        subfolders, cross-linked with \code{[[wiki-links]]} that resolve to real posts and produce
        backlinks automatically;
  \item a \textbf{learning tracker} --- per-post status, favourites, tags, external references and
        YouTube resources that play without leaving the article, and a glossary that defines a term in
        English and Vietnamese without interrupting the reading;
  \item a \textbf{knowledge graph} --- the same library seen as nodes and edges, so gaps and clusters are
        visible at a glance rather than inferred from a file tree;
  \item an \textbf{AI learning assistant} --- a chat that can read the library through explicit,
        read-only queries and \emph{propose} concrete changes: generate an article, build a learning path
        as a tree of posts, link things together, reorganise a branch. Every proposal is reviewed
        item by item before anything is written.
\end{itemize}

\subsection{The two commitments that shaped the design}

Almost every interesting decision in this project follows from two commitments made before any code was
written.

\begin{mochabox}[title={1. Local-first, not offline-tolerant}]
The knowledge base lives in a Room database on the device and nowhere else. There is no account, no sync,
no cloud copy, and Android's own Auto Backup is switched off so the operating system does not quietly
upload the database either. Consequently every feature except the assistant works with the network
switched off --- not degraded, but identical --- and the app cannot lose someone's library because a
server went away. The cost is real and accepted: no multi-device access, and backup is the user's
responsibility, which is why export/import is a first-class screen with a weekly reminder.
\end{mochabox}

\begin{mochabox}[title={2. The AI never writes to the database}]
A language model that can call \code{DELETE} is a data-loss feature. Instead the model has no database
access at all: it can ask for information through seven read-only queries, and it can \emph{propose}
changes only as a batch of structured JSON actions. The device parses that batch, validates it against
live data, shows it on a review screen where each action can be inspected and deselected, and then
applies the approved subset in a single Room transaction. Destructive actions arrive unchecked and need an
explicit confirmation. This is what makes an assistant with write ambitions safe enough to actually use.
\end{mochabox}

\subsection{Why this is not a template project}

The parts that took the most thought are the parts that do not exist in a tutorial: a single
self-referencing table that makes branches, folders and posts share one set of tree operations; a
content-backed FTS4 index kept in step by Room's own triggers; wiki-link resolution that survives
renaming a post; a model-independent action protocol with on-device validation and single-level undo; a
two-round context mechanism that lets the model query the library without ever being sent it; and a
force-directed graph renderer written against \code{Canvas} with the layout maths in framework-free Java
so it can be unit-tested. \S\ref{sec:self} is honest about what is and is not finished.
